\begin{shadequote}[l]{Richard Feynman}
    If you think you understand quantum mechanics, you don't understand quantum mechanics. 
\end{shadequote}
\ \\
Hello!

This is a compilation of lecture notes I took at Lakeside School during the 2024-25 school year on MIT OpenCourseWare's 8.04 Quantum Mechanics I course, as taught in 2016 by Prof. Barton Zweibach.\footnote{ \url{https://ocw.mit.edu/courses/8-04-quantum-physics-i-spring-2016/}} They are intended primarily as a supplement to the lectures; it can be helpful to see the words of the professor on paper, or in slightly different language. The order in which material is presented here and how it is grouped closely aligns with the structure of the lectures. 

I was advised in my independent study by Dr. Dounas-Frazer, who provided invaluable insight throughout. I also referenced Griffith's \textit{Introduction to Quantum Mechanics, 3rd Edition}\footnote{\url{https://doi.org/10.1017/9781316995433}}, which is an excellent introductory and intermediate text with plenty of exercises and problems. 

These notes are accessible to anyone who has taken multivariable calculus and has a solid understanding of classical physics. Basic linear algebra is used in Chapter 1, and we introduce inner products, Fourier transforms, Hermitian operators, and commutators. As always in physics, the more math background you possess, the easier it will be.

Quantum mechanics is one of the most fascinating, yet misunderstood, topics in modern science. In this course, you will learn the theory behind such phenomena as entanglement, tunneling, and nondeterminism. However, I want to emphasize that it \textit{is just math}. If you can solve differential equations (and if you can't, soon you will be able to), you can do quantum mechanics. 

In taking and compiling these notes, despite my best efforts at maintaining accuracy, it's very possible (and in fact quite likely, considering the length of this text) that I have made a mistake somewhere. If you catch one, please reach out to me or the Lakeside physics department to fix it! Thanks, and enjoy!

- Daniel Wang '25

\textit{May 2025}

